% Part: proof-theory
% Chapter: normalization
% Section: normalization-thm

\documentclass[../../../include/open-logic-section]{subfiles}

\begin{document}

\olfileid{pt}{nor}{thm}

\olsection{सामान्यरूपीकरण प्रमेय}

एखाद्या !!^a{proof} मध्ये
कट-खंड नसतील, तर ती
\emph{सामान्य} रूपात असते.
म्हणजे तिच्यावर क्रमपरिवर्तन
रूपांतरण किंवा न्यूनीकरण
रूपांतरण यांपैकी कोणतेही
लावता येत नाही; आणि
हेच या स्थितीचे लक्षण आहे.
एखादी !!a{proof}
\emph{सामान्यरूप करणे}
म्हणजे कट शिल्लक
राहणार नाहीत तोपर्यंत
क्रमपरिवर्तन आणि न्यूनीकरण
रूपांतरणे लावून तिला
सामान्यरूप सिद्धतेत
बदलणे. हे नेहमी शक्य
असते, हाच
\emph{सामान्यरूपीकरण प्रमेय}.

\begin{thm}\ollabel{thm:normalization}
कोणतीही~\Log{N1i}-!!{proof}
$\delta$ सामान्यरूप करता
येते; म्हणजे तिचे कट
नसलेल्या !!a{proof}~$\delta'$
मध्ये रूपांतर करता येते.
\end{thm}

\begin{proof}
  ~$\delta$ ची कट-कोटी
  आणि कट-लांबी यांवर
  विगमन करू. विगमन
  गृहीतक असे: कोणतीही
  !!{proof} जिची
  कट-कोटी कमी असेल,
  किंवा कट-कोटी समान
  पण कट-लांबी कमी असेल,
  सामान्यरूप करता येते.

  कट-लांबी~$0$
  असेल, तर कट नसतात
  आणि~$\delta$ आधीच
  सामान्यरूप असते.
  म्हणून~$\delta$ ची
  कट-कोटी आणि कट-लांबी
  दोन्ही~$>0$ आहेत असे
  समजू. सर्वात वरचा
  आणि सर्वात उजवा
  महत्तम कट-खंड निवडा.
  त्या कट-खंडाची लांबी
  $>1$ असेल, तर
  क्रमपरिवर्तन रूपांतरण
  लावा. लांबी~$1$
  असेल, तर संबंधित
  न्यूनीकरण रूपांतरण लावा.
  याने मिळणाऱ्या नव्या
  !!{proof} ची कट-कोटी
  तीच राहून कट-लांबी
  कमी होते, किंवा
  कट-कोटीच कमी होते;
  म्हणून विगमन गृहीतक
  लागू होते.
\end{proof}

वरील सिद्धतेतील
पद्धत—नेहमी सर्वात
वरचा आणि सर्वात
उजवा महत्तम कट
हाताळणे—सामान्यरूप
!!a{proof} मिळवण्यासाठी
रूपांतरणे एका ठरावीक
क्रमाने लावावी लागतात
असे सुचवते. पण आश्चर्य
म्हणजे क्रम महत्त्वाचा
नसतो. क्रमपरिवर्तन आणि
न्यूनीकरण रूपांतरणे
कोणत्याही क्रमाने
लावली तरी शेवटी
सामान्यरूप !!{proof}
मिळते.

खंड आणि कट यांच्या
व्याख्या~\Log{N2i}
साठीही सहज देता येतात;
आणि~\Log{N1i} मधील
क्रमपरिवर्तन व न्यूनीकरण
रूपांतरणे थेट~\Log{N2i}
मध्ये रूपांतरित होतात.
त्यामुळे सामान्यरूपीकरण
प्रमेय~\Log{N2i} साठीही
खरे आहे.

\begin{thm}\ollabel{thm:normalization-N2i}
कोणतीही~\Log{N2i}-!!{proof}
$\delta$ सामान्यरूप करता
येते; म्हणजे तिचे कट
नसलेल्या !!a{proof}~$\delta'$
मध्ये रूपांतर करता येते.
\end{thm}

\end{document}
